\begin{objectifs}
À la fin de ce chapitre, je sais :
\begin{itemize}[label=\textbullet]

    \item Se rappeler le vocabulaire précis des opérations. (\ii{terme},
        \ii{facteur}, \ii{somme}, etc).
    \item Additionner, soustraire, multiplier et diviser pour résoudre des
        problèmes et contrôler la vraisemblance de son résultat.
    \item Connaître le sens et les situations d'emploi de ces opérations.
    \item Enchaîner des opérations.
    \item Traduire un problème, une succession donnée d’opérations, un programme
        de calcul, en une seule expression, en faisant appel ou non à des
        parenthèses.
    \item Nommer un calcul, distinguer sommes et produits, termes et facteurs.
    \item Connaître et utiliser les priorités opératoires.
    \item Connaître et utiliser la distributivité simple sur des exemples
        numériques.
    \item Mobiliser un algorithme dans le cadre du calcul numérique.

\end{itemize}
\end{objectifs}



